% Generated draft: requires mathematical and source review.
\documentclass[11pt]{article}
\usepackage[margin=25mm]{geometry}
\usepackage{fontspec}
\setmainfont{Latin Modern Roman}
\usepackage{amsmath,amssymb,mathtools,unicode-math}
\setmathfont{Latin Modern Math}
\usepackage{enumitem,needspace,xurl,hyperref}
\setlength{\parindent}{0pt}
\setlength{\parskip}{5pt}
\setlength{\emergencystretch}{3em}
\newcommand{\N}{\mathbb{N}}
\newcommand{\Z}{\mathbb{Z}}
\newcommand{\Q}{\mathbb{Q}}
\newcommand{\R}{\mathbb{R}}
\newcommand{\C}{\mathbb{C}}
\newcommand{\F}{\mathbb{F}}
\newcommand{\A}{\mathbb{A}}
\newcommand{\PP}{\mathbb{P}}
\newcommand{\E}{\mathbb{E}}
\newcommand{\eps}{\varepsilon}
\newcommand{\dd}{\,\mathrm{d}}
\begin{document}
\section*{UnsolvedMath ID: 20003341}
\textbf{Source problem code:} catalogue entry
\subsection{Definitions and mathematical statement}
Question 46 (A. Katok) asks for constructions of periodic orbits in triangular billiards. A regular periodic orbit in a nondegenerate Euclidean triangle is a finite closed unit-speed billiard trajectory satisfying the equal-angle reflection law at relative interiors of sides, avoiding vertices, and returning with its initial direction. Iterates do not count as new primitive trajectories. The source does not limit this construction request to rational triangles.

The June 2004 source states that every acute triangle has the altitude-foot (Fagnano) orbit, while further orbits in general acute triangles and even existence in general obtuse triangles were unknown there. It cites special classes [66, 83, 35, 77, 171], and a computer-assisted result using McBilliards [84] that every triangle with all angles strictly below $100$ degrees has a periodic orbit [172]. These are dated source claims. The supplied aid separately reports a June 2026 preprint of Giovanni Forni claiming a nonconstructive existence theorem for every finite bounded polygon. It explicitly qualifies this as a very recent unrefereed v1, not an independently certified published theorem, and supplies no title or identifier. General existence, even if established, does not by itself supply the requested explicit constructions.

The aid reports two constructive companion results. First, every regular periodic orbit with an odd number $n$ of reflections has an itinerary-preserving transverse strip returning under the doubled word of length $2n$. Applied to the three-bounce Fagnano orbit, this gives a six-bounce family in every acute triangle; the central orbit may have the shorter primitive period three. All nearby paths must retain side admissibility and positive clearance from vertices.

Second, the aid describes a reflection-symmetric convex-polygon criterion requiring interior perpendicular feet; the general criterion's full side hypotheses are absent, but its stated isosceles specialization is explicit. For
\[
 T(a,h)=\operatorname{conv}\{(-a,0),(a,0),(0,h)\},\qquad a,h>0,
 \quad s=\sqrt{a^2+h^2},
\]
put $O=(0,0)$ and
\[
 P_\pm=\left(\frac{\pm ah^2}{a^2+h^2},\frac{a^2h}{a^2+h^2}\right).
\]
The feet $P_\pm$ lie in the interiors of the corresponding sloping sides. The path $O,P_+,O,P_-,O$, with its reflected directions at the two visits to $O$, is a regular primitive four-bounce trajectory, lying in a four-bounce cylinder as described by the aid. Its length and a positive vertex-clearance certificate are
\[
 \ell=\frac{4ah}{s},\qquad
 d_0=\min\left\{a,\frac{a^2}{s},\frac{h^2}{s}\right\}>0.
\]
The repeated base position has different directions, so it does not reduce the phase-space period to two.

For a rational polygon, whose angles are rational multiples of $\pi$, the source describes unfolding to a compact flat surface, with possible conical singularities, so billiard paths become straight geodesics [141]. Its quoted quadratic estimates must count periodic cylinders or families, not individual periodic phase points, of which a cylinder contains continuously many. With a fixed appropriate counting convention, the source reports constants $c,C>0$ such that
\[
 cT^2\leq N(T)\leq CT^2\qquad(T\text{ sufficiently large})
\]
[137, 138], and asymptotics $N(T)\sim c_0T^2$ for certain tables [193, 194, 52, 51], while calling the general rational-polygon asymptotic unknown at that time. The exact primitive, orientation, and unfolding-multiplicity conventions are not specified.

Finally the source announces nonuniformity near the $30$--$60$--$90$ triangle: for every $\varepsilon>0$ there is a triangle within $\varepsilon$ of it having no periodic path of length less than $1/\varepsilon$, attributed in the extraction to ``R. E. Schwatz''. A shape-space metric and length normalization are missing; the likely name spelling and the announced result require source checking. No uniform statement without those conventions is asserted here.
\subsection{Short English statement}
Construct regular periodic paths in triangles. The record combines dated existence and counting questions, explicit odd-period and isosceles construction results from a research aid, and a carefully qualified June 2026 claim of general polygonal existence. Quadratic orbit counts refer to periodic families, not the continuum of individual paths inside each cylinder.
\subsection{Sources}
\begin{itemize}
\item Local UnsolvedMath registry, permanent ID 20003341, exact source code catalogue entry; \texttt{aim-other-notes.json}, notes[147], Question 46, A. Katok.
\item American Institute of Mathematics, \emph{Emerging applications of measure rigidity}, June 2004, \url{https://aimath.org/WWN/measrigid/measrigid.pdf}. The supplied verification mentions page 19; no independent PDF inspection is asserted here.
\item Retained source labels: [66, 83, 35, 77, 171] special triangles; McBilliards [84] and the below-$100$-degree result [172]; unfolding survey [141]; quadratic bounds [137, 138]; special asymptotics [193, 194, 52, 51]. The announced nearby-triangle nonuniformity is attributed in the extraction to R. E. Schwatz. Complete bibliography entries are not supplied.
\item Supplied research attempt 1: partial\_result, problem status partially\_solved, rigor label proved, novelty confidence low. It reports the explicit constructions, counting correction and Giovanni Forni's June 2026 unrefereed v1 existence claim. This machine-generated aid is not peer reviewed; independently verify before citation. No complete preprint citation is supplied.
\end{itemize}
\end{document}
